% \iffalse meta-comment
%
% hit-thesis-backmatter.dtx 是结论、致谢、简历、成果与索引页
%
% \fi
%
% \section{结论、致谢、简历、成果与索引页}
%
%    \begin{macrocode}
%<*thesis-backmatter>
\RequirePackage{changepage}
\RequirePackage{multicol}
%    \end{macrocode}
%
% \changes{v3.2a}{2026/08/12}{\pkg{changepage} 与 \pkg{multicol} 从
%   \file{src/flow/hit-thesis-deps.dtx} 挪过来。整个类里只有下面索引那一页调它们}
% 索引那一页要比正文更宽的左右边距（\env{adjustwidth}）和双栏（\env{multicols*}），
% 这两个包只服务于它。
%
%
% \begin{environment}{conclusions}
% \changes{v3.1d}{2025/03/03}{\cs{hit@appendix@chapter*}如果不给第四个参数就不写英文目录}
%    \begin{macrocode}
\NewDocumentEnvironment { conclusions } { }
  {
    \hit@clear@page
    \hit@if@option@TF{en}%
      {\hit@appendix@chapter*{\hit@term@conclusion@heading@en}[\hit@term@conclusion@heading@en]}%
      {\hit@appendix@chapter*[\hit@term@conclusion@heading@zh]%
        {\hit@spread@by:nn{\hit@term@conclusion@heading@spread}{\hit@term@conclusion@heading@zh}}%
        [\hit@term@conclusion@heading@en]}%
  }
  { }
%    \end{macrocode}
% \end{environment}
% \begin{environment}{acknowledgements}
% \changes{v3.1d}{2025/03/03}{\cs{hit@appendix@chapter*}如果不给第四个参数就不写英文目录}
%    \begin{macrocode}
\NewDocumentEnvironment { acknowledgements } { }
  {
    \hit@clear@page
    \hit@if@option@TF{en}%
      {\hit@appendix@chapter*{\hit@term@acknowledgements@heading@en}[\hit@term@acknowledgements@heading@en]}%
      {\hit@appendix@chapter*[\hit@term@acknowledgements@heading@zh]%
        {\hit@spread@by:nn{\hit@term@acknowledgements@heading@spread}{\hit@term@acknowledgements@heading@zh}}%
        [\hit@term@acknowledgements@heading@en]}%
  }
  { }
%    \end{macrocode}
% \end{environment}
% \begin{environment}{correspondingaddr}
%    \begin{macrocode}

\NewDocumentEnvironment { correspondingaddr } { }
  { \hit@clear@page \hit@appendix@chapter * { \hit@term@resume@corresponding@address@zh } [ \hit@term@resume@corresponding@address@en ] } { }
%    \end{macrocode}
%
% 博后用到的通讯页面
% \changes{v3.0.0}{2020/05/21}{博后用到的通讯页面}
% \end{environment}
%    \begin{macrocode}


\NewDocumentEnvironment { resume } { }
  { \hit@clear@page \hit@appendix@chapter * { \hit@term@resume@heading@zh } [ \hit@term@resume@heading@en ] } { }
\NewDocumentEnvironment { publication } { }
  { \hit@clear@page \hit@appendix@chapter * { \hit@term@publications@heading@zh } [ \hit@term@publications@heading@en ] } { }


\NewDocumentEnvironment { doctorpublication } { }
  { \hit@clear@page \hit@appendix@chapter * { \hit@term@publications@doctoral@heading@zh } [ \hit@term@publications@doctoral@heading@en ] } { }
%    \end{macrocode}
%
% 此处添加博后用到的博士生发表页面
% \changes{v3.0.0}{2020/05/21}{此处添加博后用到的博士生发表页面}
% 此处中英文索引的格式设置尽量符合\PGR\ 中给出的示例的格式。此处间距常数是人工调节的。
% \changes{v1.0.10}{2018/02/19}{修改了索引的间距，使其更符合规范中的示例}
%    \begin{macrocode}
%    \end{macrocode}
%
% \changes{v3.2a}{2026/09/01}{索引页版式照规范附录 6 的官方示例改}
% 官方示例（规范附录 6，iota-hit 逐条量过）：两栏是一张假分栏的表格，
% 两列各 4230 缇、整表占满版心，栏间只有两个格的默认边距共 10.8\,bp——
% 先前的栏距 4\,em 与两侧各缩 2\,em 都无据，去掉；条目行距 1.35
% （\texttt{w:line="324" auto}），比正文的 1.25 松。
%    \begin{macrocode}
\NewDocumentEnvironment { ceindex } { }
  {
    \hit@clear@page
    \hit@appendix@chapter * { \indexname } [ \hit@term@index@heading@en ]
    \dim_set:Nn \columnsep { 10.8 bp }
    \begin { multicols* } { 2 }
    \dim_set:Nn \baselineskip { 21.01 bp }
  }
  { \end { multicols* } }
%    \end{macrocode}
%
% \changes{v3.2a}{2026/09/01}{\env{theindex} 自定义，不再走 \pkg{book} 原版}
% \emph{\env{theindex} 自己写一个。}\pkg{book} 的原版开头是
% \cs{twocolumn}\oarg{\cs{@makeschapterhead}\cs{indexname}}——\pkg{splitidx}
% 把 \cs{twocolumn} 换成“排掉可选参数”、把 \cs{@makeschapterhead} 换成
% \cs{section*}，于是每一段索引（中文一段、英文一段）开头都排一个节标题的
% 竖直空当：实测中英两段之间空出 76\,bp，章标题到首组也多出四十来 bp。
% 条目那边原版是 \cs{@idxitem}＝\cs{par}\cs{hangindent}~40\,pt，折行悬挂
% 40\,pt。官方示例（规范附录 6）里条目一律顶格、折行不悬挂，两段接着排。
%
% 所以整个环境重写：不排任何标题（章标题由 \env{ceindex} 排过了），条目
% 顶格、不悬挂，段间不加距，组间的段前由 \file{.ist} 的组标题给。两端
% 对齐关掉——一条是“词 + 页码”，拉伸没有意义。
%    \begin{macrocode}
\RenewDocumentEnvironment { theindex } { }
  {
    \dim_zero:N \parindent
    \skip_zero:N \parskip
    \raggedright
    \cs_set_protected:Npn \item { \par \dim_zero:N \hangindent }
    \cs_set_protected:Npn \subitem { \par \dim_zero:N \hangindent }
    \cs_set_protected:Npn \subsubitem { \par \dim_zero:N \hangindent }
    \cs_set_protected:Npn \indexspace { \par }
  }
  { \par }

\newlist{idxwordlist}{description}{3}
\setlist[idxwordlist,~1]{%
  itemsep=\baselineskip,
  labelindent=8em,
  font=\normalsize\bfseries,
}
\setlist[idxwordlist,~2]{%
  nosep,
  labelindent=2em,
  font=\wuhao\rm,
}
\setlist[idxwordlist,~3]{%
  nosep,
  labelindent=4em,
  font=\wuhao\rm,
}
%    \end{macrocode}
%
%    \begin{macrocode}
%</thesis-backmatter>
%    \end{macrocode}
%
% \Finale
